\paragraph{IU/Open-i chest radiograph report} (development and test only; CC BY-NC-ND 4.0).
\begin{quote}\footnotesize\raggedright
\textit{State.} EXAM: Chest radiograph\\ INDICATION: SOB\\ COMPARISON: None available\\ FINDINGS: Flattening of the bilateral hemidiaphragms. Lungs are clear. Soft tissues and bony structures unremarkable. No pneumothorax or effusion.\\ IMPRESSION: Mild hyperexpansion. No acute process.
\begin{enumerate}[label=Q\arabic*.,leftmargin=2.2em,itemsep=1pt,topsep=2pt]
  \item Is a lung nodule reported as present? \textit{(yes/no)} \textit{Answer:} no.
  \item Does the radiologist report scoliosis in this study? \textit{(yes/no)} \textit{Answer:} no.
  \item Is this a normal study with no abnormal findings? \textit{(yes/no)} \textit{Answer:} no.
\end{enumerate}
\end{quote}

\paragraph{Eurorad case report} (case 14824, \emph{Abdominal wall endometriosis: MRI findings and the “Gorgon sign”}, \url{https://www.eurorad.org/case/14824}; \textcopyright{} European Society of Radiology, CC BY-NC-SA 4.0).
\begin{quote}\footnotesize\raggedright
\textit{State.} \textit{Age:} 36\\ \textit{Sex:} female\\ \textit{Clinical history:} A 42-year-old female patient suffering from a painful “lump” at the anterior abdominal wall in the left lower quadrant. Physical examination revealed local tenderness, circumscribed hard-consistency nodule fixed at abdominal wall planes.\\ \textit{Imaging findings:} Ultrasound (Fig.1) showed a sizeable (4x3x2.5 cm) hypoechoic mass with lobulated contour, which extended ventrally from the aponeurosis of the left rectus abdominis muscle into the subcutaneous fat. The latter structures showed pre-contrast T1-hyperintensity consistent with recent blood and a fluid-fluid level. Additionally, several linear bands were seen radiating into the adjacent fat, configurating the characteristic "Gorgon sign". After intravenous gadolinium (Fig.3) both the mass and radiating strands showed marked contrast enhancement.
\begin{enumerate}[label=Q\arabic*.,leftmargin=2.2em,itemsep=1pt,topsep=2pt]
  \item Which diagnosis best explains the clinical history and imaging findings? \textit{Options:} \textbf{opt\_\allowbreak{}1: Scar endometriosis at the anterior abdominal wall}; opt\_\allowbreak{}2: Abdominal wall (incisional) hernia; opt\_\allowbreak{}3: Post-surgical seroma /haematoma; opt\_\allowbreak{}4: Soft tissue malignancy e.g. sarcoma; opt\_\allowbreak{}5: Bland (fibrotic) scar / keloid; opt\_\allowbreak{}6: Rectus muscle desmoid; opt\_\allowbreak{}7: Abdominal wall abscess; opt\_\allowbreak{}8: Other benign soft tissue tumour e.g. lipoma; opt\_\allowbreak{}9: Suture granuloma / gossypiboma. \textit{Answer:} opt\_1.
  \item Which reading worklist should this study go to? \textit{Options:} abdominal\_\allowbreak{}imaging; neuroradiology; musculoskeletal\_\allowbreak{}system; chest\_\allowbreak{}imaging; uroradiology\_\allowbreak{}genital\_\allowbreak{}male\_\allowbreak{}imaging: Kidneys, urinary tract, male genital organs; paediatric\_\allowbreak{}radiology: Children; head\_\allowbreak{}neck\_\allowbreak{}imaging; cardiovascular: Heart and vessels; \textbf{genital\_\allowbreak{}female\_\allowbreak{}imaging: Female pelvis and reproductive organs}; breast\_\allowbreak{}imaging: Breast; interventional\_\allowbreak{}radiology: Image-guided procedures. \textit{Answer:} genital\_female\_imaging.
\end{enumerate}
\end{quote}

\paragraph{MedMCQA examination question} (Apache-2.0).
\begin{quote}\footnotesize\raggedright
\textit{State.} \textit{Subject:} Physiology\\ \textit{Question:} Heart muscle, true are all except:
\begin{enumerate}[label=Q\arabic*.,leftmargin=2.2em,itemsep=1pt,topsep=2pt]
  \item Which of the options is right? \textit{Options:} \textbf{opt\_\allowbreak{}1: Has multiple nuclei}; opt\_\allowbreak{}2: Act as syncitium; opt\_\allowbreak{}3: Has branching; opt\_\allowbreak{}4: Has gap junctions. \textit{Answer:} opt\_1.
\end{enumerate}
\end{quote}

\paragraph{MedQA examination question} (CC BY 4.0).
\begin{quote}\footnotesize\raggedright
\textit{State.} \textit{Question:} A 49-year-old man comes to the physician because of a 1-week history of diarrhea and abdominal bloating. His stools are bulky, foul-smelling, and difficult to flush. Over the past 6 months, he has had recurrent dull epigastric pain that is triggered by meals and lasts for a few days. He drinks 6 to 8 beers daily. Abdominal examination shows mild epigastric tenderness with no rebound or guarding. A CT scan of the abdomen is shown. The structure indicated by the arrows is most likely lined by which of the following?
\begin{enumerate}[label=Q\arabic*.,leftmargin=2.2em,itemsep=1pt,topsep=2pt]
  \item Which option correctly answers the question? \textit{Options:} opt\_\allowbreak{}1: Pyogenic membrane; opt\_\allowbreak{}2: Columnar mucinous epithelium; \textbf{opt\_\allowbreak{}3: Granulation tissue}; opt\_\allowbreak{}4: Simple ductal epithelium. \textit{Answer:} opt\_3.
\end{enumerate}
\end{quote}

\paragraph{MMLU examination question} (development and test only; MIT).
\begin{quote}\footnotesize\raggedright
\textit{State.} \textit{Subject:} clinical knowledge\\ \textit{Question:} Which of the following molecules does not form part of DNA?
\begin{enumerate}[label=Q\arabic*.,leftmargin=2.2em,itemsep=1pt,topsep=2pt]
  \item Which of the options is right? \textit{Options:} opt\_\allowbreak{}1: Deoxyribose; \textbf{opt\_\allowbreak{}2: Amino acid}; opt\_\allowbreak{}3: Pyrimidine; opt\_\allowbreak{}4: Purine. \textit{Answer:} opt\_2.
\end{enumerate}
\end{quote}

\paragraph{PubMedQA research question} (development and test only; MIT).
\begin{quote}\footnotesize\raggedright
\textit{State.} \textit{Research question:} Does screening or surveillance for primary hepatocellular carcinoma with ultrasonography improve the prognosis of patients?\\ \textit{Abstract:} The purpose of this paper is to evaluate the efficacy of ultrasonographic screening for primary hepatocellular carcinoma. A total of 680 eligible cases were classified into three groups (surveillance, opportunistic, and symptomatic groups) according to their initial exposure. We used survival time, tumor morphology, and T staging as prognostic outcomes. The outcomes of screened/unscreened and sur veillance/nonsur veillance were compared with the use of the logistic regression model. The adjusted odds ratios for the screened group versus the unscreened group, with 1-, 2-, and 3-year survival time being used as outcomes, were 0.33 (95\% confidence interval [CI], 0.21-0.52), 0.33 (95\% CI, 0.21-0.53), and 0.37 (95\% CI, 0.23-0.61), respectively. The adjusted odds ratios for surveillance versus nonsurveillance were 0.58 (95\% CI, 0.35-0.97), 0.45 (95\% CI, 0.27-0.74), and 0.44 (95\% CI, 0.26-0.73). The odds ratios were even smaller when tumor morphology or T stage was taken as the main outcome. All these results were statistically significant. There were significant gradient relationships between prognostic outcomes and extent of screening history.
\begin{enumerate}[label=Q\arabic*.,leftmargin=2.2em,itemsep=1pt,topsep=2pt]
  \item What do the study's results say about the question? \textit{Options:} \textbf{yes: Yes, the evidence supports it}; no: No, the evidence does not support it; maybe: Maybe: the evidence is mixed or inconclusive. \textit{Answer:} yes.
\end{enumerate}
\end{quote}

\paragraph{MedXpertQA examination question} (development and test only; MIT).
\begin{quote}\footnotesize\raggedright
\textit{State.} \textit{Question:} A 54-year-old woman with a history of diabetes and peripheral vascular disease presents with swelling and aching pain in her calves that worsens after prolonged sitting. On examination, her posterior tibialis and dorsalis pedis pulses are graded 2/4 bilaterally. Both legs are symmetrically enlarged and non-pitting, with a negative Homan’s sign on both sides. General tenderness is noted from the ankles to the knees without localized areas of pain. To directly address her condition, what is the most appropriate next step in treatment?
\begin{enumerate}[label=Q\arabic*.,leftmargin=2.2em,itemsep=1pt,topsep=2pt]
  \item Choose the correct answer. \textit{Options:} opt\_\allowbreak{}1: Venous pump activation exercises; opt\_\allowbreak{}2: Pedal pump technique; opt\_\allowbreak{}3: Thoracic inlet release technique; opt\_\allowbreak{}4: Pelvic diaphragm release technique; opt\_\allowbreak{}5: Lower extremity muscle energy technique; opt\_\allowbreak{}6: Effleurage of the bilateral legs; opt\_\allowbreak{}7: Lymphatic drainage massage; opt\_\allowbreak{}8: Myofascial release of the gastrocnemius; opt\_\allowbreak{}9: Compression stockings application; \textbf{opt\_\allowbreak{}10: Indirect thoracoabdominal diaphragm release}. \textit{Answer:} opt\_10.
\end{enumerate}
\end{quote}

\paragraph{CT-RATE chest CT report} (CC BY-NC-SA 4.0; access-controlled, report text not reproduced).
\begin{quote}\footnotesize\raggedright
\textit{State.} A CT-RATE chest CT report with clinical information, technique, findings and impression.
\begin{enumerate}[label=Q\arabic*.,leftmargin=2.2em,itemsep=1pt,topsep=2pt]
  \item Does this report describe a pleural effusion? \textit{(yes/no; three such questions per report)}
  \item Is this a normal study with no abnormal findings? \textit{(yes/no)}
  \item Which of these findings is described in the report? \textit{(one finding present and three absent)}
\end{enumerate}
Answers follow the dataset's classifier labels.
\end{quote}

\paragraph{Teacher-labeled questions} (states from CheXpert Plus, ReXGradient-160K and CT-RATE reports and from clinical indications, including Eurorad case presentations).
\begin{quote}\footnotesize\raggedright
\textit{Report question.} State: a radiology report. What follow-up does this report call for? \textit{Options:} No further imaging recommended; Follow-up imaging within 3 months; Follow-up imaging in more than 3 months; Further evaluation with a different imaging test or procedure; Clinical correlation or non-imaging workup only.\\
\textit{Order question.} State: the clinical indication only. Which imaging study is most appropriate first for this clinical question? \textit{Options:} Chest radiograph; CT head without contrast; CT pulmonary angiography (PE protocol); CT chest; CT abdomen and pelvis with IV contrast; Ultrasound abdomen (incl. RUQ); Ultrasound pelvis / transvaginal; MRI brain; MRI spine; Radiograph of the painful bone or joint; MRI of a joint or limb; Doppler ultrasound of vessels; Mammography and/or breast ultrasound; No imaging indicated.\\
The label is the option that both teachers ranked first; the training target is the mean of the two teachers' distributions.
\end{quote}
